% Generic entity registry and rendering helpers.
\newcounter{characterCount}
\newcounter{spaceCount}
\newcounter{clueCount}
\newcounter{mysteryCount}
\newif\ifMMRegistrationPass
\MMRegistrationPassfalse

\makeatletter
\newcommand{\MMDeclareNamedEntity}[4]{%
  \@ifundefined{mm@#1@#2}{%
    \expandafter\gdef\csname mm@#1@#2\endcsname{#3}%
    \stepcounter{#4}%
  }{%
    \PackageError{murder-mystery}{Duplicate #1 identifier `#2'}%
      {Each entity identifier must be declared only once in its own content file.}%
  }%
}

\newcommand{\MMNamedEntity}[2]{%
  \@ifundefined{mm@#1@#2}{%
    \PackageError{murder-mystery}{Unknown #1 identifier `#2'}%
      {Declare this identifier in its entity file and add that file to the matching index.tex.}%
    \textbf{[unknown #1: #2]}%
  }{%
    \csname mm@#1@#2\endcsname%
  }%
}
\makeatother

% Raw names are useful in headings; references apply the styles configured in config/game.tex.
\newcommand{\CharacterName}[1]{\MMNamedEntity{character}{#1}}
\newcommand{\SpaceName}[1]{\MMNamedEntity{space}{#1}}
\newcommand{\ClueName}[1]{\MMNamedEntity{clue}{#1}}

\newcommand{\Character}[1]{\CharacterReferenceStyle{\CharacterName{#1}}}
\newcommand{\Space}[1]{\SpaceReferenceStyle{\SpaceName{#1}}}
\newcommand{\Clue}[1]{\ClueReferenceStyle{\ClueName{#1}}}

\newcommand{\GameTitlePage}{%
  \begin{titlepage}
    \centering
    \vspace*{0.20\textheight}
    {\color{wine}\rule{0.7\textwidth}{1.2pt}\par}
    \vspace{1.5em}
    {\Huge\bfseries\color{ink}\GameTitle\par}
    \vspace{0.65em}
    {\Large\itshape \GameTagline\par}
    \vspace{1.5em}
    {\color{wine}\rule{0.7\textwidth}{1.2pt}\par}
    \vfill
    \begin{minipage}{0.72\textwidth}
      \centering\color{ink}\GamePremise
    \end{minipage}\par
    \vfill
    \GameDate\quad\textbullet\quad\GameLocation
  \end{titlepage}%
}

% Each sheet command is its entity's single source of truth. During the silent
% registration pass it records the ID and name; during normal rendering it typesets the body.
\newcommand{\CharacterSheet}[4]{%
  \ifMMRegistrationPass
    % A character without a name is referenced by its title.
    \if\relax\detokenize{#3}\relax
      \MMDeclareNamedEntity{character}{#1}{#2}{characterCount}%
    \else
      \MMDeclareNamedEntity{character}{#1}{#3}{characterCount}%
    \fi
  \else
    \subsubsection{#2: \Character{#1}}
    \label{character:#1}%
    #4%
  \fi
}

\newcommand{\SpaceSheet}[3]{%
  \ifMMRegistrationPass
    \MMDeclareNamedEntity{space}{#1}{#2}{spaceCount}%
  \else
    \subsubsection{\Space{#1}}
    \label{space:#1}%
    #3%
  \fi
}

\newcommand{\ClueSheet}[3]{%
  \ifMMRegistrationPass
    \MMDeclareNamedEntity{clue}{#1}{#2}{clueCount}%
  \else
    \subsubsection{\Clue{#1}}
    \label{clue:#1}%
    #3%
  \fi
}

\newcommand{\Mystery}[2]{%
  \ifMMRegistrationPass
    \stepcounter{mysteryCount}%
  \else
    \subsubsection{#1}%
    #2%
  \fi
}

\newcommand{\RelationshipSummary}{%
  \subsection{Entity summary}
  There are \arabic{characterCount} characters.\par
  There are \arabic{clueCount} clues.\par
  There are \arabic{mysteryCount} misteries.\par
}
